% Section: The core: Base
\section{The core: Base}
\label{sec:base}

The Base library (\code{src/core/Base}) provides what every other part of \CAP{} builds on: the base class of named
objects and its reporting of messages, exceptions, a tree of named objects of any type, a registry of classes for
reflection, the configuration of tasks, the binary file format, and the reading and writing of XML formatted data. Its classes
depend only on the C++ standard library, with the exception of few helper optional  functions used  to load and export data from ROOT files.

% ---------------------------------------------------------------------------------------------------------------
\subsection{Object model}
\label{sec:objects}

\paragraph{Objects and messages.}
Most classes of \CAP{} derive from \code{Object}, which gives them a name, a title, a class name and a version
index. Each object has a minimum report level, one of the severities \code{TRACE}, \code{DEBUG}, \code{INFO},
\code{WARNING}, \code{ERROR} and \code{FATAL}; a message is printed only if its severity is at least that level, with
a header naming the class, the object and the function that issued it. The report level of each task is set in its
configuration, so that the messages of one part of a workflow can be followed without those of the others. All
exceptions of \CAP{} derive from \code{Exception}, itself a \code{std::exception}, and record a message, the function
that threw them and their class (file, configuration, histogram, mathematical domain, and so on).

\paragraph{Folders.}
A \code{Folder} is a container of named objects of any type and of sub-folders, forming a tree. The folder owns its
objects: copying a folder copies them, and removing or replacing an object deletes it. Objects are reached by name or
by path through the sub-folders, \code{get<Histogram>("PtPt/pairs/n2")}, and a folder can be navigated like a file
system, with \code{cd()} and relative or absolute paths. Typed access works through the base classes of the objects,
so that a histogram is found as a \code{Histogram}, a \code{DataPoints} or an \code{Object}. Two folders of the same
structure can be added or subtracted entry by entry, for every type that has a sum, to combine the results of
several jobs; sub-samples are combined by a dedicated task, which averages them and takes their spread as the
uncertainty (Section~\ref{sec:subsample}). Folders are the unit of input and output of the analyses: each job writes the
histograms of its tasks to a folder, with one sub-folder per task and one per combination of filters
(Section~\ref{sec:folders-analysis}), and saves it to a file (Section~\ref{sec:binary}).

\paragraph{Reflection.}
Earlier versions of \CAP{} relied on the dictionaries of ROOT to create objects from their class names, which is
how the tasks of a workflow are instantiated from an XML file. \CAP{}~8 has its own run-time registry of classes
(\code{Reflection.hpp}), written in standard C++17. A class is described once in C++: its bases, its constructors,
its data members and properties, and its member functions, overloads and inherited functions included. The registry
then finds classes by name and lists them by base class, creates objects with any of their constructors, reads and
writes their properties, and calls their member functions by name:
\begin{verbatim}[language=C++]
const ClassInfo * h = Registry::instance().find("Histogram");
Value g    = Registry::instance().find("MersenneTwister")->create({ 12345 });
Value hist = h->create({ "h", "Gaussian numbers", 40, -4.0, 4.0 });
for (int k=0; k<10000; k++) hist.call("fill", { g.call("gauss", { 0.5, 1.0 }) });
hist.call("nBins", { 0 });   // a member function of Partition, a base of Histogram
\end{verbatim}
All the task classes are registered, one line each, so that workflows create their tasks by class name. The same
registry serves the command-line interpreter, CapShell (Section~\ref{sec:capshell}). Its descriptions are written by
hand today; they are meant to be generated by the static reflection of C++26 once the compilers used by \CAP{}
support it.

\paragraph{Tasks.}
The work of a job is done by a tree of tasks (\code{Task}). A task goes through a fixed life cycle: its default
configuration is defined (\code{setDefaultConfiguration()}), the configuration is read (\code{configure()}), the task
is initialised once, executed repeatedly, typically once per event, and finalised. A task owns its subtasks and
passes each step of the cycle on to them, after doing its own part; an event iterator, for instance, executes a
generator and then the analyses for each event. Tasks are created from their class names through the registry,
and their parameters, subtasks included, are read from XML (Section~\ref{sec:xml}).

% ---------------------------------------------------------------------------------------------------------------
\subsection{Binary data model}
\label{sec:binary}

\CAP{}~8 stores its objects in its own binary format, independent of ROOT. The format was designed to be compact,
portable between platforms and compilers, robust against corrupted or truncated files, and readable by later
versions of the code.

\paragraph{Records.}
A file is a sequence of records, one per object. Each record is
\begin{center}
  \small
  \begin{tabular}{@{}lll@{}}
    \toprule
    Field & Size & Content \\
    \midrule
    tag & 4 bytes & \code{CAP} followed by one letter identifying the class \\
    version & 1 byte & format version of the payload \\
    length & varint & number of bytes of the payload \\
    payload & length & the name of the object, then the data of its class \\
    checksum & 4 bytes & CRC-32 of the payload \\
    \bottomrule
  \end{tabular}
\end{center}
Integers are written as variable-length integers (LEB128: seven bits per byte, one byte below 128), signed integers
after a zigzag mapping that keeps small negative numbers short, floating-point numbers as their IEEE~754 bit patterns,
and every multi-byte value in little-endian order, byte by byte, so that files are the same on every platform.
Strings are written as their length followed by their bytes. Arrays of doubles are written densely or sparsely
(only the non-zero values, each with the gap in index since the previous one), whichever is smaller, and are restored
exactly. The CRC-32 checksum (the polynomial of zip and png) detects corrupted data, and every read is checked against
the bounds of the record, so that a damaged file is refused with an exception rather than read wrongly.
Table~\ref{tab:tags} lists the classes that are stored this way.

\begin{table}[t]
  \centering
  \caption{Classes stored in the CAP binary format, with their record tags.}
  \label{tab:tags}
  \small
  \begin{tabular}{@{}llll@{}}
    \toprule
    Tag & Class & Tag & Class \\
    \midrule
    \code{CAPO} & \code{Folder} (holds the records of its entries) & \code{CAPF} & \code{ScalarField} \\
    \code{CAPH} & \code{Histogram} & \code{CAPM} & \code{Matrix} \\
    \code{CAPQ} & \code{Profile} & \code{CAPV} & \code{VectorND} \\
    \code{CAPD} & \code{DataPointsVector} & \code{CAP2}, \code{CAP3} & \code{Vector2D}, \code{Vector3D} \\
    \code{CAPK} & \code{DataPointColumns} & \code{CAPL} & \code{VectorLorentz} \\
    \code{CAPP}, \code{CAPT}, \code{CAPN} & \code{DataPoint2D}, \code{3D}, \code{ND} & \code{CAPR}, \code{CAPI} & \code{Parameter}, \code{LongParameter} \\
    \code{CAPB} & \code{Table2D} & & \\
    \bottomrule
  \end{tabular}
\end{table}

\paragraph{Histograms.}
The payload of a histogram holds its name and title, then its rank and, for each axis, the number of bins, the range
and the title; then the number of non-empty bins and, for each of them, the gap in global index since the previous
non-empty bin, the number of entries, the average value and its uncertainty; then the labels of the bins and the fits
attached to the histogram. Empty bins are not written at all, and a filled bin takes about ten bytes, so that the
large and sparse histograms of the pair analyses stay small on disk. The content of the histogram is restored
exactly.

\paragraph{Functions and fits.}
A function is stored as part of the fits attached to data points, a histogram or another function. A fit record holds
the name under which the class of the fitted function is registered (Section~\ref{sec:functions}), its number of
parameters and its rank, the name, value, type and bounds of each parameter, the label and result of the fit, its
range, and its regions. When the file is read, the function is recreated by name through the registry and given its
parameters, so that any registered class, a composite registered under a name included, is read back. A fit whose
class is not registered by the program reading the file is skipped, and the number of fits skipped is returned to
the caller.
\todo{decide whether functions should also be stored as objects of their own in folders, not only as fits.}

\paragraph{Folders and files.}
A folder is stored as one \code{CAPO} record whose payload holds the records of its entries, sub-folders included,
so that a whole tree of results is one record. When a file is read, each record is rebuilt as an object of its
class through a registry of types, in which the Math library registers its classes when it is loaded; programs add
their own with \code{Folder::registerType<T>()}. A \code{CapFile} holds one or several folders. It can be read eagerly,
or lazily: the tree of folders is then built at once, with the names and types of all the objects, and each object is
decoded only when it is first used, which makes it quick to open a large file and look at a few of its histograms.
The checksums of a lazily read file are verified when it is opened, so a corrupted file is still refused at once.

\paragraph{Versions.}
Each class writes its records with its current format version and reads every version up to that one: histograms,
for instance, are at version~4 and still read the earlier versions. Records of classes a program does not know are
skipped, so that one file can hold objects of many classes. Additions that older code can ignore are placed at the end
of a block: the regions of a fit (Section~\ref{sec:fitting}) are written after its range, and only when the fit has
some, so that files written with regions are still read, without them, by earlier versions of \CAP{}~8.

% ---------------------------------------------------------------------------------------------------------------
\subsection{Configuration in XML}
\label{sec:xml}

\paragraph{Properties and configurations.}
The parameters of the tasks are properties: named values whose names are paths of components separated by colons,
such as \code{PtPt:HISTOGRAM:pt\_nbins}. Values are stored as text and read as booleans, integers, floating-point
numbers or strings; a property keeps the type of its default value, so that a value of the wrong type is detected when
a configuration is read. A key matches a property of the same name or one whose name ends with it, so that a task reads
its parameters with keys relative to its own name. A \code{Configuration} adds numbered keys, lists of values such as
the particle filters of a task (\code{PARTICLE\_FILTER\_0}, \code{PARTICLE\_FILTER\_1}, \dots).

\paragraph{XML.}
The configuration of a task maps to an XML element: the task is a \code{TASK} element whose name, class and severity
are attributes, each component of a key but the last becomes a \code{GROUP} element, the last one a \code{PARAMETER}
element holding the value and its type, and the subtasks are nested \code{TASK} elements. The XML classes of Base
(\code{XmlDocument}, \code{XmlElement}, \code{XmlReader}, \code{XmlWriter}) need nothing but the standard library. The
reader is a non-validating parser that loads no external content and does not recurse, so that malformed or deeply
nested input is reported with its line and column instead of being read wrongly; the writer streams large documents
without building them in memory first.

\paragraph{Templates.}
There is no fixed dictionary of parameters. Each class defines its own in its function
\code{setDefaultConfiguration()}, which lists its parameters with their default values and types. Run with the option \code{-{}-xml-templates}, the \code{CAP} program writes this
template for every registered task class, one XML file per class in \code{DB/Templates}. The workflow editor \code{RunCAP.py}
(Section~\ref{sec:run}) builds its task dialogs from these templates, and the values read from a workflow file are
checked against them; a parameter that is not in the template of its class is reported.

\paragraph{Workflow files.}
A workflow file has a \code{CAP} root element with one top-level \code{TASK} element per stage of the analysis (event
processing, derived histograms, balance functions), marked by a \code{stage} attribute, each with the job settings
(output and input files) and its tree of tasks. A \code{RUNCAP} element at the end holds what the editor needs to rebuild
its own view: the definitions of the event and particle filters, the generators and particle databases used, and the
histogram settings shared between tasks. An excerpt:
\begin{widetext}
\begin{verbatim}[language=XML]
<CAP version="1">
  <TASK name="Workflow" class="CAP::Task" severity="INFO" stage="EventProcessing">
    <GROUP name="JOB">
      <PARAMETER name="OUTPUT_FILE" type="string">RunAna.cap</PARAMETER>
    </GROUP>
    <TASK name="RunAnalysis" class="CAP::EventIterator" severity="INFO">
      <TASK name="PythiaEventGenerator" class="CAP::PythiaEventGenerator"> ... </TASK>
      <TASK name="ParticlePairAnalyzer" class="CAP::ParticlePairAnalyzer"> ... </TASK>
    </TASK>
  </TASK>
  <TASK name="DerivedHistos" class="CAP::Task" stage="DerivedHistos"> ... </TASK>
  <RUNCAP version="1">
    <OBJECT kind="PARTICLE_FILTER" name="PiP" title="#pi^{+}">
      <CONDITION subtype="PDG" form="value" minimum="211" maximum="211" />
    </OBJECT>
    ...
  </RUNCAP>
</CAP>
\end{verbatim}
\end{widetext}